% ======================================================================
% Appendix D -- Course Plans and Rubrics
% ======================================================================

\chapter{Course Plans and Rubrics}
\label{app:course}

The book supports one-semester courses, and shorter courses of seven or eight weeks, for graduate and upper-level undergraduate students in computer science, data science, information systems and business analytics, and MBA technology electives. Instructors choose the \textbf{build track} (design plus implementation and deployment), the \textbf{design track} (no programming), or both, with mixed teams.

\section{A 14-week plan}

\begingroup
\small
\begin{xltabular}{\textwidth}{@{}c>{\raggedright\arraybackslash}p{3.8cm}>{\raggedright\arraybackslash}X@{}}
\caption{A 14-week course.}\label{tab:course14}\\
\toprule
Week & Chapters & Work due \\
\midrule
\endfirsthead
\toprule
Week & Chapters & Work due \\
\midrule
\endhead
1 & 1 From building to designing & Lab 1; Appendix A setup \\
2 & 2 The designer's toolkit & Lab 2; decomposition exercise \\
3 & 3 The business of AI solutions & Lab 3: value case \\
4 & 4 Policy assistant & Lab 4: first deployment \\
5 & 5 Support triage & Lab 5: MCP tool design \\
6 & 6 Invoice to payment & Lab 6: checkpoint \\
7 & 7 Incident response & Lab 7; capstone teams formed \\
8 & Midterm review; capstone proposals & Capstone brief and process map \\
9 & 8 Fraud review & Lab 8: evaluation report \\
10 & 9 Disruption rebooking & Lab 9; capstone solution graphs \\
11 & 10 KYC and AML & Lab 10: controls evidence \\
12 & 11 Forecasting, or 12 Loan memo (choose one; the other as reading) & Lab 11 or 12; capstone trade-offs \\
13 & 13 Capstone & Capstone build and evaluation \\
14 & Presentations & Memo, deployment, evaluation report \\
\bottomrule
\end{xltabular}
\endgroup

\section{A 10-week plan}

Cover Chapters 1--3 in weeks 1--2; choose four cases (we suggest 4, 5, 6 and 8, which cover deployment, tool design, human-in-the-loop and evaluation) for weeks 3--6; Chapter~13 in week 7; capstone work in weeks 8--9; presentations in week 10. Assign the remaining cases as reading with their review questions.

\section{An 8-week plan}

A short course, such as a half-semester module or an executive-education program, should keep the whole design method\index{design method} and use fewer cases. Chapters~1--3 are essential, because every case depends on them. Chapters~4--6 come next: together they exercise each building block once---retrieval in the policy assistant, MCP tool design in support triage, and approval gates\index{approval gate} in invoice-to-payment---and all three start from a published sample, which is easier in a short course than designing from a blank page. The capstone becomes a smaller final project (Table~\ref{tab:course8}).

\begingroup
\small
\begin{xltabular}{\textwidth}{@{}c>{\raggedright\arraybackslash}p{3.8cm}>{\raggedright\arraybackslash}X@{}}
\caption{An 8-week course.}\label{tab:course8}\\
\toprule
Week & Chapters & Work due \\
\midrule
\endfirsthead
\toprule
Week & Chapters & Work due \\
\midrule
\endhead
1 & 1 From building to designing & Lab 1; Appendix A setup; free-trial sign-up \\
2 & 2 The designer's toolkit & Lab 2: solution graph and trade-off table for returns \\
3 & 3 The business of AI solutions & Lab 3: value case and one-page memo \\
4 & 4 Policy assistant & Lab 4: first deployment \\
5 & 5 Support triage & Lab 5: MCP tool design; project teams formed \\
6 & 6 Invoice to payment & Lab 6: checkpoint; project proposal \\
7 & One elective case (8, 9 or 7), or project work & Project solution graph and trade-off table \\
8 & Presentations; Chapter~13 rollout sections as reading & Memo, deployment (build track), evaluation summary \\
\bottomrule
\end{xltabular}
\endgroup

\paragraph{Choosing the elective in week 7.} Pick the case that fits the audience. Chapter~8 (fraud review) suits analytics and MBA students: asymmetric errors\index{asymmetric error cost} and evaluation. Chapter~9 (disruption rebooking) suits computer science students: handoffs between agents and surge planning. Chapter~7 (incident response) suits an IT or operations audience: event triggers and safe actions. Assign the other two as reading with their review questions.

\paragraph{The final project.} Chapters~10--12 (KYC and AML, forecasting and replenishment, commercial loan underwriting) form the project menu. Each team takes one as its starting point, or brings a process from its own organization, and produces the four deliverables of every case: a solution graph, a trade-off table\index{trade-off table} with a decision record\index{decision record}, a one-page memo\index{one-page memo}, and, on the build track, a working deployment. Teams form in week~5 so they can work on the project alongside Chapters~5 and~6.

\paragraph{A 7-week variant.} Drop week~7. Teams start the final project after Chapter~6 and present in week~7; the elective cases become reading.

\paragraph{Cloud credit.} In a course this short, students sign up for the free trial in week~1; the 90-day credit covers the whole course (Section~\ref{sec:app-account}).

\paragraph{Assessment.} For the short course, a suggested weighting is labs 30\%, case assignments (Chapters~4--6) 30\%, final project 30\%, and participation in design reviews 10\%.

\section{Assessment}

A suggested weighting: labs 30\%, case assignments 25\%, capstone 35\%, participation in design reviews 10\%. Case assignments are the end-of-chapter portfolio piece: solution graph, trade-off table, deployment (build track) and memo.

\section{Rubrics}

\begingroup
\scriptsize \setlength{\tabcolsep}{4pt}
\begin{xltabular}{\textwidth}{@{}>{\raggedright\arraybackslash}p{2.2cm}>{\raggedright\arraybackslash}X>{\raggedright\arraybackslash}X>{\raggedright\arraybackslash}X@{}}
\caption{Rubric for a case assignment or the capstone. Each criterion is scored 1--4.}\label{tab:rubric}\\
\toprule
Criterion (weight) & 4: exemplary & 2: developing & 1: missing \\
\midrule
\endfirsthead
\toprule
Criterion (weight) & 4: exemplary & 2: developing & 1: missing \\
\midrule
\endhead
Problem and baseline (10\%) & Process map with measured or well-sourced times, costs and error points; SUAO complete & Process described; numbers asserted without sources & No baseline \\
Solution graphs (20\%) & Two or more alternatives, well formed (W1--W8); every block choice justified & One design, or alternatives that differ trivially; some rule violations & Diagram without the notation or without alternatives \\
Trade-offs and decision (15\%) & Constraints first; money rows; sensitivity; failure walk-through changes the design; decision record with reversal triggers & Scores without reasons or sensitivity & Choice asserted \\
Business case (20\%) & Floor case from measured inputs; the case's key assumption identified and tested by the pilot plan; one-page memo the sponsor could act on & Value estimated but not separated into floor and upside; memo too long or too technical & No value case \\
Build and deployment (15\%) & Deployed; tools follow T1--T9; identity and secrets handled; cleanup done & Runs locally only, or tools API-shaped & Not built \\
Evaluation (15\%) & Metrics tied to KPIs and risks; thresholds set in advance; adversarial and failure tests; results honestly reported & Some metrics, mostly answer quality & No evaluation \\
Communication (5\%) & Clear, concise, plain language; graphs readable & Understandable with effort & Unclear \\
\bottomrule
\end{xltabular}
\endgroup

For the design track, replace ``Build and deployment'' with ``Build plan'' (node cards for every agent, server and guardrail; deployment target decision with reasons) and keep the other criteria.

\section{Running design reviews in class}

Design reviews are the most valuable class activity this book supports. A format that works: each team presents its solution graph for five minutes; a second team, assigned in advance, plays the reviewer, applying W1--W8, asking for the trade-off table, and choosing one node for a failure walk-through\index{failure walk-through} on the spot; the instructor plays the sponsor and asks the questions of Chapter~\ref{ch:business}: what changes, what is it worth, what could go wrong, how will we know. Rotate roles weekly. Students learn as much from reviewing as from presenting.

\section{The free design coach}

The companion design coach---available as a Claude artifact and a custom GPT---reviews a student's solution graph (as a list of nodes and edges or a picture) and trade-off table against the checklists in Appendix~C, and points back to the relevant section of the book. It holds chapter summaries and rubrics, not the book's text, and it does not grade. Instructors may require students to run their design through the coach before a review, so that class time is spent on judgment rather than on checklist items.
